\section{Summary and Discussion.}

We first give an overview of various ideas that have influenced us
and to which we have not yet referred directly. Subsequently we
proceed to a summary explaining what are the essential new points
we have made in this paper. We finish with a short description of
planned future work.

\subsection{Connections to other work.}

At the conceptual level, this paper has been influenced by the
work of many others, in particular people in our own subfield of
lattice field theory, even if this was not mentioned at any point
in our presentation until now. While it is impossible to avoid
omissions, we try here to draw attention to what stands out in our
memory at this point.

That an extra dimension might be needed for a local string dual of
the field theoretical description clearly is an idea with many
sources. The relationship between this extra dimension and a space
of eigenvalues has been emphasized by many, for example in papers
by A. Jevicki \cite{jevi}. The essential point is that for a
hermitian or unitary matrix the eigenvalues are restricted to a
one dimensional line and repel each other; as such they make up
dynamically, at infinite $N$, a continuous one dimensional degree
of freedom. In the context of matrix models this produces an
analogue of the extra dimension we have seen in higher dimensional
setups.

Using the lattice to make progress on the question of a possible
string dual of nonabelian YM theory has preoccupied a sizable
group of researchers. This new wave of lattice work on the large $N$
limit has been started by M. Teper, and in various collaborations
he has addressed many of the questions we address
here \cite{teper}. In particular, the question of the cross-over
has been studied in various contexts \cite{teper-cross-over}. There
is an interesting issue raised there suggesting a qualitative
difference between three dimensional and four dimensional YM gauge
theories with respect to observable dependent large $N$
transitions. The immediate implication is that it would be
interesting to repeat our work here in three Euclidean dimensions;
it would be surprising to see three dimensions so different in
this context from both two and four dimensions.

J. Kuti \cite{kuti} and collaborators have studied in great detail the cross-over
from field theoretical to string-like behavior and have
shown how various
string states reorganize themselves as one goes from short to long
QCD strings. The philosophy of matching various field theories
with known or suspected string-like excitations onto an effective
string theory at long distances underlies much of this work. This
work has not yet been extended to large values of $N$, but this is
recognized as a possible valid avenue for further steps in this
program.

We have been influenced significantly by papers by R. Brower
\cite{brower}. On the one hand side he and collaborators have
provided additional examples in which one can see the extra
dimension being connected to standard scale decompositions of four
dimensional processes of traditional interest. This interpretation
is very natural when one thinks in terms of a warp factor in the
context of the famous AdS/CFT correspondence in the conformal and
related cases. It provides a most convincing resolution of the old
puzzle of how a string theory can embody hard, field theoretic
ultraviolet behavior.

It has also been stressed by Brower that one should be more
ambitious than merely connecting onto an effective long distance
string theory, and, that the extra dimension could provide the
means to construct a complete stringy dual to YM, valid at all
distances, albeit too difficult for actual calculation at short
distances. It is with this in mind that although our strings are
fat, we maintain the zigzag symmetry, which would be a potential
clue about the infinitely thin string dual. Also, with this in
mind, we insisted on finding observables that have a continuum
field theoretical semiclassical expansion at small sizes, where
the effects of short distance fluctuations reveal themselves.

The effective string description has been studied in various
guises also by others \cite{other}, in addition to Polchinski
and Strominger; the approaches focused both on
open Dirichlet strings and on closed string wrapping round a large
compact direction of space-time.

Another related topic is the issue of ``Casimir scaling'', in the
context of the dependence of the string tension on the
representation of the source and sink the $SU(N)$ string connects.
It is known that ``Casimir scaling'' is exact when applied to
charges in the fundamental and in the adjoint representation
at $N=\infty$,
so long as we take $N\to\infty$ first and let the minimal area
spanned by the loop increase only after that~\cite{adlerneu}.
We wish to mention here a recent paper~\cite{negele} in the context
of ``Casimir scaling'' at finite $N$ where the
Wilson loop eigenvalue distribution has been studied on the
lattice for $SU(2)$.

Last but not least we wish to mention the old idea of
preconfinement \cite{amativ16}, in which one pictures a jet
producing process in QCD as evolving by first creating (in the
planar limit) finite energy clusters that already are color
neutral and only subsequently hadronize. This preparation for
hadronization is similar to the early signals of closing the gap
in the eigenvalue spectrum, before the density becomes uniform up
to corrections exponential in the minimal area spanned by the
loop, where one has full confinement. As mentioned previously, one
can have a large $N$ transition without ever making it all the way
to confinement.

\subsection{Main new ideas in this work.}

Our main change in vantage was a focus on the universal nature of
large $N$ transition which we hope to turn into a calculation that
can take us through the cross-over which connects the regime in
which field theory is the most convenient description to the
regime in which string theory is assumed to eventually become a
convenient description.

We see the beginnings of such a calculational scheme and feel that
were it successful, it would provide a concrete result on which to
build in the search for a more elegant and fundamental dual
representation of Yang Mills theory.

We found ourselves naturally led into something that looks like an
extra dimension and one that is intimately related to scaling at
that.


\subsection{Plans for the future.}

The first problem for the future is a more thorough investigation
of the critical regime of the large $N$ transitions and the
numerical identification of the right universality class. In
parallel an effort will be made to make the semiclassical
calculation based on instantons at the short distance end
concrete. The double scaling limits in the larger universality
class would need to be studied in detail; assuming that our guess
for this universality class is right, it should be possible to
construct those double scaling limits exactly.

The connection onto a convenient effective string model at the
strong coupling end is another direction of future activity; here,
the main tool would be numerical at first, and we expect to draw
from the experience of the many other workers on similar problems.
Our main focus will be on keeping $N$ large so the string coupling
constant can be set to zero as is often done even for $N=3$.

Also, we feel it might be useful to try to face the question
whether an effective string theory for planar QCD is indeed not
different from an effective string theory in a model which has
nothing to do with non-abelian gauge theory. For example, on the
face of it one would expect the zigzag symmetry to play no role at
leading order in the effective string theory, but maybe this is
not a correct assumption. An analogy is the chiral Lagrangian
describing pions in QCD, which indeed could come from a
fundamental theory that has no gauge fields. But, the complete
decoupling at all scales of all the mesons, well known to occur at
infinite $N$, is something the effective Lagrangian, by itself,
does not incorporate in a structural way; rather, many of its free
parameters need to be set to zero.
